% !TeX program = lualatex
% !TeX encoding = utf8
% !TeX spellcheck = uk_UA
% !TeX root = ../main.tex

%% --------------------------------------------------------
%% Chapter file: Fluid.tex
%% --------------------------------------------------------

\graphicspath{{\currfilebase/Pictures}}

\Opensolutionfile{answer}[\currfilebase/\currfilebase-Answers]
\Writetofile{answer}{\protect\section*{Розділ~\ref*{sec:mhd}. \nameref*{sec:mhd}}}
\chapter{Гідродинамічний та магнітогідродинамічний опис}\label{sec:mhd}
\label{\currfilebase}\hypertarget{mhd}{}
%% --------------------------------------------------------

До цього моменту ми розглядали окремі частинки. Але в~багатьох задачах не потрібно знати
траєкторію кожної з~них: достатньо знати усереднені поля густини, швидкості, тиску та
електромагнітного поля.


\section{Гідродинамічний опис не є синонімом частих зіткнень}


Зіткнення полегшують встановлення локальної рівноваги та ізотропного тиску. Проте
великомасштабні низькочастотні рухи плазми можуть описуватися рідинними рівняннями і~за
рідких зіткнень, якщо швидкі кінетичні ступені вільності не визначають основну динаміку.
Саме тому в~плазмі співіснують два описи: кінетичний, у~термінах функцій розподілу, та
МГД, у~термінах макроскопічних полів.

Для однорідинної МГД типові умови мають вигляд
\begin{equation*}
L\gg r_{Li},\qquad \omega\ll\omega_{ci},
\end{equation*}
але цього недостатньо саме по собі: потрібні також відповідне замикання рівнянь та, для
найпростішої ізотропної моделі, достатньо швидка ізотропізація тиску. Якщо ці умови
порушуються, треба повертатися до дворідинної або кінетичної теорії.


%% --------------------------------------------------------
\section{Дворідинний опис}\index{Дворідинний опис|textbf}
%% --------------------------------------------------------


\begin{Theory}[Нагадаємо з~гідродинаміки]
Для звичайної рідини з~густиною $\rho$ і~швидкістю $\vect v$ маємо рівняння неперервності
$\partial_t\rho+\nabla\cdot(\rho\vect v)=0$ та~рівняння Ейлера
\[\rho\Bigl(\frac{\partial\vect v}{\partial t}+(\vect v\cdot\nabla)\vect v\Bigr)=-\nabla p+\vect f,\]
де $\vect f$~--- об'ємна сила; в'язкість додає член Нав'є--Стокса $\eta\nabla^2\vect v$. Зверніть увагу
на \emph{проблему замикання}: рівняння для $\rho$ містить $\vect v$, рівняння для $\vect v$~--- тиск $p$,
рівняння для $p$~--- потік тепла і~т.~д. Ланцюжок обривають моделлю стану (наприклад, $p\propto\rho^\gamma$).
Нижче рівняння для електронної та іонної \enquote{рідин} читаються як \enquote{Ейлер + сила Лоренца + тертя $\vect R$}.
\end{Theory}

Найприродніший гідродинамічний крок~--- не відразу зливати електрони та іони в~одну
рідину, а~описувати їх окремо. Для кожного компонента вводимо концентрацію $n_s$ і~середню
швидкість $\vect v_s$. Рівняння неперервності мають вигляд
\begin{equation*}
 \frac{\partial n_s}{\partial t}
 +\nabla\cdot(n_s\vect v_s)=0.
\end{equation*}
Рівняння руху містять силу Лоренца, градієнт тиску та силу зіткнень:
\begin{equation*}
 m_s n_s\frac{d\vect v_s}{dt}
 =
 q_s n_s\left(\vect E+\frac1c\vect v_s\times\vect B\right)
 -\nabla p_s+\vect R_s.
\end{equation*}
Сума $\sum_s\vect R_s=0$: зіткнення передають імпульс між компонентами, але не створюють
його. Тут ці рівняння записано за аналогією з~гідродинамікою; їхній систематичний вивід як
моментів кінетичного рівняння (і~звідки береться тензор тиску замість скалярного $p_s$) наведено
в~розд.~\ref{sec:kinetic-moments}.

Для квазінейтральної водневої плазми $n_e\simeq n_i=n$ та $m_e\ll m_i$. Важливо побачити,
що однорідинна МГД не вводиться незалежно від дворідинної моделі: її основні рівняння
можна отримати безпосередньо, додаючи та віднімаючи рівняння для електронів та іонів.


\section{Від дворідинної моделі до однорідинної МГД}\index{Магнітна гідродинаміка (МГД)|textbf}


Для однозарядної плазми рівняння руху мають вигляд
\begin{equation*}
m_i n_i\frac{D_i\vect v_i}{Dt}
=
en_i\left(\vect E+\frac1c\vect v_i\times\vect B\right)
-\nabla p_i+\vect R_{ie},
\end{equation*}
\begin{equation*}
m_e n_e\frac{D_e\vect v_e}{Dt}
=-en_e\left(\vect E+\frac1c\vect v_e\times\vect B\right)
-\nabla p_e+\vect R_{ei},
\end{equation*}
де $\vect R_{ei}=-\vect R_{ie}$ виражає взаємну передачу імпульсу під час зіткнень.


\subsection*{Сума рівнянь: рівняння руху плазми як цілого}


Введемо масову густину та масову швидкість плазми:
\begin{equation*}
\rho_m=m_i n_i+m_e n_e,
\qquad
\vect v=
\frac{m_i n_i\vect v_i+m_e n_e\vect v_e}{\rho_m}.
\end{equation*}
Для $m_e\ll m_i$ маємо $\rho_m\simeq m_i n_i$ і~$\vect v\simeq\vect v_i$. Додаючи рівняння
руху електронів та іонів, внутрішні сили зіткнень взаємно скорочуються. У~квазінейтральній
плазмі $ n_e\simeq n_i$ скорочується також сумарна електрична сила:
\begin{equation*}
en_i\vect E-en_e\vect E\simeq0.
\end{equation*}
Магнітні сили, навпаки, не скорочуються, оскільки електрони та іони рухаються з~різними
швидкостями. Використовуючи
\begin{equation*}
\vect j=e n_i\vect v_i-e n_e\vect v_e,
\end{equation*}
їхня сума набуває вигляду
\begin{equation*}
\frac{e n_i}{c}\vect v_i\times\vect B
-
\frac{e n_e}{c}\vect v_e\times\vect B
=
\frac1c\vect j\times\vect B.
\end{equation*}
Якщо нехтувати внеском відносного руху електронів та іонів у~потік механічного імпульсу
(для звичайної низькочастотної МГД це мала поправка), дістаємо
\begin{equation}\label{eq:one-fluid-momentum-from-two-fluid}
\boxed{
\rho_m\frac{D\vect v}{Dt}
=
-\nabla(p_i+p_e)
+
\frac1c\vect j\times\vect B.
}
\end{equation}
Позначивши $p=p_i+p_e$, одержуємо рівняння руху плазми як однієї (однорідинної) рідини.

Точно так само, додаючи рівняння неперервності двох компонент, отримуємо рівняння
неперервності для єдиної рідини:
\begin{equation}\label{eq:one-fluid-continuity-from-two-fluid}
\boxed{
\frac{\partial\rho_m}{\partial t}
+
\nabla\cdot(\rho_m\vect v)=0.
}
\end{equation}


\subsection*{Електронне рівняння: закон Ома}\index{Ома закон!узагальнений|textbf}


Тепер скористаємося рівнянням для електронів. Для квазінейтральної плазми покладемо
$n_e\simeq n_i=n$ та, для простоти, $\vect v\simeq\vect v_i$. Тоді
\begin{equation*}
\vect v_e=\vect v-\frac{\vect j}{en}.
\end{equation*}
Підставимо це у~рівняння руху електронів і~покладемо $m_e\to0$ (інерція електронів нехтовна):
\[0=-en\Bigl(\vect E+\frac{\vect v_e\times\vect B}c\Bigr)-\nabla p_e+\vect R_{ei},\qquad
\vect R_{ei}=\frac{en}{\sigma}\,\vect j.\]
Зіткнення електронів з~іонами дають силу тертя $\vect R_{ei}=en\vect j/\sigma$ (знак відповідає гальмуванню відносного руху).
Підстановка $\vect v_e=\vect v-\vect j/en$ дає
$-en\bigl(\vect E+\vect v\times\vect B/c\bigr)+\vect j\times\vect B/c-\nabla p_e+en\vect j/\sigma=0$, і~після ділення на $en$:
\begin{equation}\label{eq:generalized-ohm-from-two-fluid}
\tcbhighmath{
    \vect E+\frac1c\vect v\times\vect B
    =
    \frac{\vect j}{\sigma}
    +
    \frac{\vect j\times\vect B}{enc}
    -
    \frac{\nabla p_e}{en}.
}
\end{equation}
Тут перший доданок праворуч описує опір зіткненням, другий є \emphx{холлівським}\index{Холла член|textbf} доданком,
а~третій виникає через градієнт електронного тиску. Холлівський доданок важливий на масштабах порядку іонної інерційної довжини $c/\omega_{pi}$: на менших масштабах іони \enquote{відклеюються}
від поля, а~електрони ще ні (токамак, $n=10^{14}~\mathrm{см}^{-3}$: $c/\omega_{pi}\approx2$--$3$~см; магнітосфера, $n\sim1~\mathrm{см}^{-3}$: $\approx230$~км).
Саме ці доданки порушують вмороженість і~уможливлюють перез'єднання силових ліній (розд.~\ref{sec:reconnection}). Якщо нехтувати всіма цими поправками,
отримуємо закон ідеальної МГД:
\begin{equation}\label{eq:ideal-mhd-ohm}
\tcbhighmath{
\vect E+\frac1c\vect v\times\vect B=0.
}
\end{equation}

Таким чином, перехід має чітку фізичну структуру:
\begin{equation*}
\tcbhighmath{
\begin{array}{c}
\text{сума рівнянь двох компонент}
\\[-1mm]
\Downarrow
\\[-1mm]
\text{рівняння руху однієї рідини}
\\[1mm]
\text{електронне рівняння}
\\[-1mm]
\Downarrow
\\[-1mm]
\text{узагальнений закон Ома}
\\[-1mm]
\Downarrow
\\[-1mm]
\text{ідеальна МГД}
\end{array}}
\end{equation*}

Саме тому МГД слід розглядати як довгохвильову низькочастотну границю більш повного
дворідинного опису, а~не як незалежну систему рівнянь.


\section{Основне рівняння магнітної гідродинаміки}
%% --------------------------------------------------------


Перейдемо до макроскопічного опису плазми, розглядаючи її як суцільне провідне середовище
з~масовою густиною $\rho_m(\vect r,t)$, гідродинамічною швидкістю $\vect v(\vect r,t)$ та
газокінетичним тиском $p(\vect r,t)$.

Виведемо рівняння руху такого середовища, послідовно розглянувши сили, що діють на його
елемент. Ми отримаємо те саме рівняння~\eqref{eq:one-fluid-momentum-from-two-fluid}, але
тепер~--- з~балансу сил на~елементі об'єму, що робить наочним зміст кожного доданка.


\subsection*{Сила газокінетичного тиску}


Виділимо в~плазмі елементарний об'єм $dV=dx\,dy\,dz$ (рис.~\ref{pic:volume-element}).
Перепад тиску на протилежних гранях уздовж осі $x$ створює силу

% ~~~~~~~~~~~~~~~~~~~~~~~~~~~~~~~~~~~~~~~~~~~~~~~~~~~~~~~~
\begin{equation*}
    df_x
    =
    [p(x,y,z)-p(x+dx,y,z)]\,dy\,dz
    =
    -\frac{\partial p}{\partial x}\,dV.
\end{equation*}
% ~~~~~~~~~~~~~~~~~~~~~~~~~~~~~~~~~~~~~~~~~~~~~~~~~~~~~~~~

Сумуючи внески вздовж трьох осей координат, дістаємо повну силу тиску, що діє на елемент
об'єму:

% ~~~~~~~~~~~~~~~~~~~~~~~~~~~~~~~~~~~~~~~~~~~~~~~~~~~~~~~~
\begin{equation*}
    d\vect f_p=-\nabla p\,dV.
\end{equation*}
% ~~~~~~~~~~~~~~~~~~~~~~~~~~~~~~~~~~~~~~~~~~~~~~~~~~~~~~~~

\begin{figure}[ht]
    \centering
    \localinput{volume-element.tikz}
    % \includegraphics[width=0.4\linewidth]{volume-element}
    \caption{Елемент об'єму в~плазмі, на який діють сили газокінетичного тиску}
    \label{pic:volume-element}
\end{figure}


\subsection*{Рівняння руху та сила Лоренца}


Маса виділеного елемента дорівнює

% ~~~~~~~~~~~~~~~~~~~~~~~~~~~~~~~~~~~~~~~~~~~~~~~~~~~~~~~~
\begin{equation*}
    dm=\rho_m\,dV.
\end{equation*}
% ~~~~~~~~~~~~~~~~~~~~~~~~~~~~~~~~~~~~~~~~~~~~~~~~~~~~~~~~

Тому за другим законом Ньютона

% ~~~~~~~~~~~~~~~~~~~~~~~~~~~~~~~~~~~~~~~~~~~~~~~~~~~~~~~~
\begin{equation*}
    \rho_m\,dV\,\frac{d\vect v}{dt}
    =
    d\vect f_p+d\vect f_{\mathrm{em}}.
\end{equation*}
% ~~~~~~~~~~~~~~~~~~~~~~~~~~~~~~~~~~~~~~~~~~~~~~~~~~~~~~~~

Тут

% ~~~~~~~~~~~~~~~~~~~~~~~~~~~~~~~~~~~~~~~~~~~~~~~~~~~~~~~~
\begin{equation*}
    \frac{d}{dt}
    =
    \frac{\partial}{\partial t}
    +(\vect v\cdot\nabla)
\end{equation*}
% ~~~~~~~~~~~~~~~~~~~~~~~~~~~~~~~~~~~~~~~~~~~~~~~~~~~~~~~~
--- субстанціальна похідна, яка враховує як явну зміну поля швидкості з~часом, так і~його
зміну під час руху разом із плазмою.

У~квазінейтральній плазмі густина нескомпенсованого заряду мала, тому безпосередньою
електричною силою $\rho_q\vect E$ у~макроскопічному рівнянні руху нехтуємо. Магнітна
складова сили Лоренца при цьому залишається істотною:

% ~~~~~~~~~~~~~~~~~~~~~~~~~~~~~~~~~~~~~~~~~~~~~~~~~~~~~~~~
\begin{equation*}
    d\vect f_{\mathrm{em}}
    =
    \frac{1}{c}\,\vect j\times\vect B\,dV.
\end{equation*}
% ~~~~~~~~~~~~~~~~~~~~~~~~~~~~~~~~~~~~~~~~~~~~~~~~~~~~~~~~

Отже,

% ~~~~~~~~~~~~~~~~~~~~~~~~~~~~~~~~~~~~~~~~~~~~~~~~~~~~~~~~
\begin{equation*}
    \rho_m\frac{d\vect v}{dt}
    =
    -\nabla p
    +
    \frac{1}{c}\,\vect j\times\vect B.
\end{equation*}
% ~~~~~~~~~~~~~~~~~~~~~~~~~~~~~~~~~~~~~~~~~~~~~~~~~~~~~~~~

Це рівняння має простий фізичний зміст: рух плазми визначається градієнтом газокінетичного
тиску та магнітною силою Лоренца.


\subsection*{Виключення густини струму за законом Ампера}


Розглянемо повільні макроскопічні процеси, для яких струмом зміщення в~законі Ампера можна
знехтувати. Тоді

% ~~~~~~~~~~~~~~~~~~~~~~~~~~~~~~~~~~~~~~~~~~~~~~~~~~~~~~~~
\begin{equation*}
    \vect j=\frac{c}{4\pi}\,\rot\vect B.
\end{equation*}
% ~~~~~~~~~~~~~~~~~~~~~~~~~~~~~~~~~~~~~~~~~~~~~~~~~~~~~~~~

Це зручно, бо в~плазмі струми зазвичай невідомі, а~поле~--- те, що ми задаємо котушками.
Підставляючи цей вираз у~рівняння руху, отримуємо \emphx{основне рівняння магнітної
гідродинаміки} (МГД):

% ~~~~~~~~~~~~~~~~~~~~~~~~~~~~~~~~~~~~~~~~~~~~~~~~~~~~~~~~
\begin{equation}\label{eq:MHD-equation}
    \tcbhighmath{
    \rho_m\frac{d\vect v}{dt}
    =
    -\nabla p
    +
    \frac{1}{4\pi}\,
    (\rot\vect B)\times\vect B}.
\end{equation}
% ~~~~~~~~~~~~~~~~~~~~~~~~~~~~~~~~~~~~~~~~~~~~~~~~~~~~~~~~

Отже, на відміну від звичайної гідродинаміки, на рух плазми безпосередньо впливає власне
магнітне поле~--- через створювану ним силу Лоренца.

Скориставшись векторною тотожністю $(\rot\vect B)\times\vect B=(\vect B\cdot\nabla)\vect
B-\nabla(B^2/2)$, перепишемо його. (Тотожність випливає з~$\nabla(B^2/2)=(\vect B\cdot\nabla)\vect B+\vect B\times(\rot\vect B)$,
або в~компонентах: $[(\rot\vect B)\times\vect B]_i=\varepsilon_{ijk}\varepsilon_{jlm}\partial_lB_m\,B_k=B_k\partial_kB_i-B_k\partial_iB_k$.)
% ~~~~~~~~~~~~~~~~~~~~~~~~~~~~~~~~~~~~~~~~~~~~~~~~~~~~~~~~
\begin{equation}\label{eq:MHD-pressure-tension}
	\rho_m\frac{d\vect v}{dt}=-\nabla\left(p+\frac{B^2}{8\pi}\right)+\frac1{4\pi}\,(\vect B\cdot\nabla)\vect B.
\end{equation}
% ~~~~~~~~~~~~~~~~~~~~~~~~~~~~~~~~~~~~~~~~~~~~~~~~~~~~~~~~~~~~~~~~
Тут два фізично різні доданки. Перший~--- \emph{магнітний тиск} $B^2/8\pi$, що додається
до газового. Другий~--- \emph{натяг силових ліній}: вздовж викривленої лінії з~радіусом
кривизни $R_c$ діє сила $B^2/(4\pi R_c)$, напрямлена до центра кривизни, ніби лінії~---
розтягнуті гумові шнури.


\subsection*{Замикання системи та межі застосовності}\index{Замикання системи рівнянь|textbf}

Рівняння~\eqref{eq:MHD-equation} саме по собі не утворює замкненої системи. Для визначення
густини, швидкості, тиску та магнітного поля необхідні додаткові рівняння.

Зокрема, до рівняння руху додають:

\begin{enumerate}

    \item рівняння неперервності
    % ~~~~~~~~~~~~~~~~~~~~~~~~~~~~~~~~~~~~~~~~~~~~~~~~~~~~~~~~
    \begin{equation*}
        \frac{\partial\rho_m}{\partial t}
        +
        \nabla\cdot(\rho_m\vect v)
        =0;
    \end{equation*}
    % ~~~~~~~~~~~~~~~~~~~~~~~~~~~~~~~~~~~~~~~~~~~~~~~~~~~~~~~~

    \item рівняння стану, наприклад адіабатичне
    % ~~~~~~~~~~~~~~~~~~~~~~~~~~~~~~~~~~~~~~~~~~~~~~~~~~~~~~~~
    \begin{equation*}
        \frac{d}{dt}\bigl(p\rho_m^{-\gamma}\bigr)=0
        \quad(\text{тобто }p\rho_m^{-\gamma}=\mathrm{const}\text{ вздовж руху елемента плазми});
    \end{equation*}
    % ~~~~~~~~~~~~~~~~~~~~~~~~~~~~~~~~~~~~~~~~~~~~~~~~~~~~~~~~

	\item закон Ома, який пов'язує $\vect E$, $\vect B$, $\vect v$ та~$\vect j$. У~простій
	  резистивній МГД електричне поле в~системі відліку, що рухається разом із~плазмою,
	  $\vect E+\vect v\times\vect B/c$, задає омічну складову струму:
    % ~~~~~~~~~~~~~~~~~~~~~~~~~~~~~~~~~~~~~~~~~~~~~~~~~~~~~~~~
    \begin{equation}\label{eq:Ohm_law_plasma}
        \vect j=\sigma\left(\vect E+\frac{1}{c}\,\vect v\times\vect B\right).
    \end{equation}
    % ~~~~~~~~~~~~~~~~~~~~~~~~~~~~~~~~~~~~~~~~~~~~~~~~~~~~~~~~

Це \emph{резистивний} закон Ома для однорідинної МГД; $\sigma$ тут~--- скаляр, тобто $\sigma_\parallel$:
поперек поля опір повністю іонізованої плазми лише приблизно вдвічі більший ($\eta_\perp\approx2\eta_\parallel$,
Спітцер), а~не в~$(\omega_c/\nu)^2$ разів. Точніший дворідинний закон
Ома~\eqref{eq:generalized-ohm-from-two-fluid} містить також холлівський доданок
$\vect j\times\vect B/enc$ і~градієнт електронного тиску $-\nabla p_e/en$, а~якщо не нехтувати
інерцією електронів~--- ще й~$-(m_e/e)\,d\vect v_e/dt$. У~простій резистивній МГД всі ці
поправки вважаються малими.

    % ~~~~~~~~~~~~~~~~~~~~~~~~~~~~~~~~~~~~~~~~~~~~~~~~~~~~~~~~
    \item рівняння Максвелла, зокрема
    % ~~~~~~~~~~~~~~~~~~~~~~~~~~~~~~~~~~~~~~~~~~~~~~~~~~~~~~~~
    \begin{equation*}
        \nabla\cdot\vect B=0,
        \qquad
        \frac{\partial\vect B}{\partial t}
        =
        -c\,\rot\vect E.
    \end{equation*}
    % ~~~~~~~~~~~~~~~~~~~~~~~~~~~~~~~~~~~~~~~~~~~~~~~~~~~~~~~~

\end{enumerate}

МГД-наближення призначене для великомасштабних ($L\gg r_{Li}$) та повільних
($\omega\ll\omega_{ci}$) процесів, для яких кінетичні ефекти, пов'язані з~індивідуальним
рухом частинок, не визначають основну динаміку плазми. Пояснення \enquote{на пальцях}: \enquote{рідина} має сенс, коли частинка за час зміни поля встигає багато разів
обернутися ($\omega\ll\omega_{ci}$), а~ларморівське коло мале порівняно з~масштабом неоднорідності ($L\gg r_{Li}$): тоді частинку \enquote{прив'язано}
до силової лінії, і~поперек $\vect B$ плазма поводиться як суцільне середовище навіть без зіткнень. Уздовж $\vect B$ ця прив'язка не працює.

%% --------------------------------------------------------
\section{Альвенівська швидкість і~число Альвена}
%% --------------------------------------------------------


Рівняння МГД показує, що магнітне поле має не лише тиск, а~й~натяг. Це означає, що
викривлення силових ліній може породжувати хвилю: магнітний натяг $B^2/4\pi$ відіграє
роль натягу струни, а~$\rho_m$~--- роль лінійної інерції. Така \emph{альвенівська хвиля}\index{Альвена хвиля}
біжить уздовж $\vect B_0$ зі швидкістю
\begin{equation}\label{eq:alfven-wave}
    \tcbhighmath{v_A=\frac{B_0}{\sqrt{4\pi\rho_m}}}
\end{equation}
(повне виведення з~лінеаризованих рівнянь МГД та інші моди~--- у~розд.~\ref{sec:mhd-waves}).

Альвенівська швидкість $v_A$~--- один із~головних масштабів швидкості в~МГД: з~нею порівнюють
швидкість течії плазми, а~час $L/v_A$ є характерним часом магнітних процесів на масштабі $L$.
Значення $v_A$ для токамака, сонячної корони та сонячного вітру наведено в~табл.~\ref{tab:alfven-numbers}.
Відношення
\begin{equation*}
M_A=\frac{v}{v_A}
\end{equation*}
називають числом Альвена. За $M_A\ll1$ магнітні сили здатні істотно впливати на рух, тоді
як за $M_A\gg1$ інерційний рух домінує над альвенівським відновлювальним механізмом. Це
порівняння є локальним і~не є самостійним критерієм застосовності МГД.

Альвенівська хвиля~--- лише одна з~хвиль замагніченої плазми. Зовнішнє поле $\vect B_0$ робить
реакцію плазми анізотропною: рух уздовж поля та поперек
нього фізично різний. У~низькочастотній границі $\omega\ll\omega_{ci}$ це дає МГД-хвилі
(альвенівську, швидку й~повільну магнітозвукові), а~при вищих частотах потрібен тензорний
опис холодної замагніченої плазми (моди $R$, $L$, $O$, $X$, резонанси та відсікання).
Обидва випадки систематично розглянуто в~розд.~\ref{sec:mhd-waves}.


%% --------------------------------------------------------
\section{МГД-рівновага та магнітні поверхні}\index{Рівновага!магнітогідродинамічна|textbf}\index{Магнітні поверхні|textbf}\label{sec:mhd-equilibrium}
%% --------------------------------------------------------



Розглянемо тепер стаціонарну плазму, в~якій $\vect v=0$. Рівняння руху перетворюється на
умову рівноваги:

% ~~~~~~~~~~~~~~~~~~~~~~~~~~~~~~~~~~~~~~~~~~~~~~~~~~~~~~~~
\begin{equation}\label{eq:equilibrium}
    \nabla p
    =
    \frac{1}{c}\,\vect j\times\vect B.
\end{equation}
% ~~~~~~~~~~~~~~~~~~~~~~~~~~~~~~~~~~~~~~~~~~~~~~~~~~~~~~~~

Це рівняння показує, що в~рівновазі градієнт тиску компенсується магнітною силою Лоренца.

Корисно мати ще один еквівалентний опис магнітної сили. Використовуючи закон Ампера та
$\nabla\cdot\vect B=0$, можна записати
\begin{equation}\label{eq:maxwell-stress}
    f_i=\partial_j\sigma_{ij},
    \qquad
    \sigma_{ij}=
    \frac{1}{4\pi}
    \left(B_iB_j-\frac{B^2}{2}\delta_{ij}\right).
\end{equation}
У~цьому тензорі перший доданок відповідає натягу вздовж поля, а~другий --- магнітному
тиску. Такий запис особливо зручний для задач про сили на плазмові шнури та магнітні
поверхні.

Скалярно помножимо~\eqref{eq:equilibrium} на $\vect B$. Оскільки $\vect B\cdot(\vect
j\times\vect B)=0$, маємо

% ~~~~~~~~~~~~~~~~~~~~~~~~~~~~~~~~~~~~~~~~~~~~~~~~~~~~~~~~
\begin{equation*}
    \vect B\cdot\nabla p=0.
\end{equation*}
% ~~~~~~~~~~~~~~~~~~~~~~~~~~~~~~~~~~~~~~~~~~~~~~~~~~~~~~~~

Отже, магнітні силові лінії лежать на поверхнях сталого тиску:

% ~~~~~~~~~~~~~~~~~~~~~~~~~~~~~~~~~~~~~~~~~~~~~~~~~~~~~~~~
\begin{equation*}
    p=\mathrm{const}.
\end{equation*}
% ~~~~~~~~~~~~~~~~~~~~~~~~~~~~~~~~~~~~~~~~~~~~~~~~~~~~~~~~

Аналогічно, скалярний добуток~\eqref{eq:equilibrium} на $\vect j$ дає

% ~~~~~~~~~~~~~~~~~~~~~~~~~~~~~~~~~~~~~~~~~~~~~~~~~~~~~~~~
\begin{equation*}
    \vect j\cdot\nabla p=0,
\end{equation*}
% ~~~~~~~~~~~~~~~~~~~~~~~~~~~~~~~~~~~~~~~~~~~~~~~~~~~~~~~~
%
тобто густина струму також лежить у~дотичній площині до поверхонь сталого тиску.

У~тороїдальних магнітних системах такі поверхні можуть утворювати вкладені одна в~одну
магнітні поверхні. Їхня геометрія є основою опису рівноваги плазми в~токамаках
і~стелараторах.

З~\eqref{eq:equilibrium} можна також знайти поперечну складову густини струму. Векторно
помноживши це рівняння на $\vect B$, дістаємо

% ~~~~~~~~~~~~~~~~~~~~~~~~~~~~~~~~~~~~~~~~~~~~~~~~~~~~~~~~
\begin{equation*}
    \vect j_\perp
    =
    \frac{c}{B^2}\,
    \vect B\times\nabla p.
\end{equation*}
% ~~~~~~~~~~~~~~~~~~~~~~~~~~~~~~~~~~~~~~~~~~~~~~~~~~~~~~~~

Це саме діамагнітний струм~\eqref{eq:diamagnetic-current}, який детальніше розглянуто в~розд.~\ref{sec:diamagnetism}.


\subsection*{Параметр плазми \texorpdfstring{$\beta$}{beta}}\index{Бета плазми|textbf}


Важливим безрозмірним параметром плазми є відношення газокінетичного тиску до магнітного:

% ~~~~~~~~~~~~~~~~~~~~~~~~~~~~~~~~~~~~~~~~~~~~~~~~~~~~~~~~
\begin{equation}\label{eq:beta}
    \tcbhighmath{
    \beta
    =
    \frac{p}{B^2/(8\pi)}
    =
    \frac{8\pi p}{B^2}.
    }
\end{equation}
% ~~~~~~~~~~~~~~~~~~~~~~~~~~~~~~~~~~~~~~~~~~~~~~~~~~~~~~~~

Параметр $\beta$ показує, наскільки сильно тиск плазми порівняно з~магнітним полем впливає
на рівновагу. Він \emph{не} є критерієм застосовності МГД: ідеальна МГД справедлива за будь-якого
$\beta$ (за умов $L\gg r_{Li}$, $\omega\ll\omega_{ci}$), тоді як $\beta$ лише вказує, наскільки
сильно плазма деформує поле.

\begin{itemize}

    \item При $\beta\ll1$ магнітний тиск значно перевищує газокінетичний, тому плазма
      відносно слабо впливає на конфігурацію магнітного поля.

    \item При $\beta\sim1$ газокінетичний та магнітний тиски порівнянні, тому взаємний
      вплив плазми та магнітного поля стає істотним.

\end{itemize}

Таким чином, $\beta$ характеризує співвідношення двох основних механізмів, що визначають
рівновагу плазми~--- її внутрішнього тиску та магнітного утримання. Для заданого
магнітного поля велике значення $\beta$ відповідає великому тиску плазми, однак досягнення
такого стану визначається також умовами магнітного утримання та стійкістю плазми. Для порівняння: токамак $\beta\sim3\%$, сферичний токамак до $\sim40\%$, магнітосфера Землі $\beta\ll1$,
сонячна корона $\sim10^{-3}$--$10^{-2}$, фотосфера $\gtrsim1$, сонячний вітер $\sim1$.

Тиск $p$ визначається термодинамічним станом плазми; для класичної ідеальної плазми він
задається формулою~\eqref{eq:ideal-pressure}.

\begin{figure}[ht]
    \centering
    \localinput{beta.tikz}
    \caption{Фізичний зміст параметра $\beta$: порівняння газокінетичного та магнітного
    тисків. Баланс тисків вимагає, щоб $\beta$, обчислене за зовнішнім (вакуумним) полем, не
    перевищувало~1; локально $\beta$ може бути й~більшим (напр., у~конфігурації з~оберненим
    полем). Практичну межу задає стійкість: у~токамаках $\beta\sim$ кілька відсотків}
    \label{pic:plasma-beta}
\end{figure}




\section{Енергія електромагнітного поля та плазми}


Теорема Пойнтінга в~системі СГС має вигляд
\begin{equation}\label{eq:poynting-theorem-plasma}
\frac{\partial}{\partial t}\left(\frac{E^2+B^2}{8\pi}\right)
+\nabla\cdot\left(\frac{c}{4\pi}\vect E\times\vect B\right)
=-\vect j\cdot\vect E.
\end{equation}
Ліва частина описує зміну густини енергії поля та її потік, права~--- локальну передачу
енергії від поля до заряджених частинок. Якщо $\vect j\cdot\vect E>0$, поле віддає плазмі
енергію.

У~МГД енергетичний баланс пов'язує рух середовища, магнітне поле та нагрів. Для ідеальної
МГД (без дисипації) з~адіабатичним замиканням повна густина енергії
$\mathcal E=\rho_mv^2/2+p/(\gamma-1)+B^2/8\pi$ (кінетична, теплова й~магнітна) зберігається:
\begin{equation}\label{eq:mhd-energy}
\frac{\partial\mathcal E}{\partial t}
+\nabla\cdot\left[\Bigl(\frac{\rho_mv^2}2+\frac{\gamma p}{\gamma-1}\Bigr)\vect v
+\frac{c}{4\pi}\,\vect E\times\vect B\right]=0,
\qquad
\frac{c}{4\pi}\vect E\times\vect B=\frac{B^2\vect v-(\vect v\cdot\vect B)\vect B}{4\pi},
\end{equation}
де останній вираз~--- потік Пойнтінга в~ідеальній МГД ($c\vect E=-\vect v\times\vect B$). З~урахуванням
скінченної провідності в~балансі магнітної й~кінетичної енергій з'являється дисипативний член
$-j^2/\sigma$ (джоулеве тепло); якщо врахувати й~теплову енергію, повна енергія знову
зберігається: джоулеве тепло лише переводить магнітну та кінетичну енергії в~теплову. Кожен доданок має чітке значення: робота
$\vect v\cdot(\vect j\times\vect B)/c$ перекачує енергію між рухом і~магнітним полем; поле
вморожене, тому розтягування силових ліній збільшує $B^2/8\pi$ коштом кінетичної енергії,
а~обмін між кінетичною й~тепловою енергіями відбувається через роботу стиснення $p\,\nabla\cdot\vect v$
(у~рівнянні для теплової енергії) та~$-\vect v\cdot\nabla p$ (у~рівнянні для кінетичної); у~сумі вони дають
потік $\nabla\cdot(p\vect v)$. Це дає й~фізичний
спосіб перевіряти наближення: якщо відкидання певного члена порушує очікуваний
енергетичний баланс, воно не може бути зроблене безумовно.


%% --------------------------------------------------------
\section{Діамагнетизм плазми}\index{Діамагнетизм!плазми|textbf}\label{sec:diamagnetism}
%% --------------------------------------------------------


Кожна заряджена частинка плазми у~магнітному полі має орбітальний магнітний
момент~\eqref{eq:mu}, напрямлений проти поля. Якщо в~одиниці об'єму є $n_s$ частинок сорту
$s$, то намагніченість
% ~~~~~~~~~~~~~~~~~~~~~~~~~~~~~~~~~~~~~~~~~~~~~~~~~~~~~~~~
\begin{equation*}
	\vect M=-\sum_sn_s\mu_s\frac{\vect B}{B}=-\sum_sn_s\frac{T_s}{B}\frac{\vect B}{B}=-\frac p{B^2}\,\vect B,
\end{equation*}
% ~~~~~~~~~~~~~~~~~~~~~~~~~~~~~~~~~~~~~~~~~~~~~~~~~~~~~~~~
де враховано $\langle mv_\perp^2/2\rangle=T$ (два ступені вільності поперек
поля) і~$p=\sum_sn_sT_s$. Намагніченість напрямлена проти поля, тобто
плазма \emph{діамагнітна}: її магнітна сприйнятливість
($M=\kappa H$, $H\approx B$) від'ємна, $\kappa=-p/B^2<0$, тому вона виштовхується з~області
сильного поля.

Діамагнетизм виявляється через струми. У~неоднорідній плазмі ($p$ змінюється в~просторі)
вони дорівнюють $\vect j=c\,\nabla\times\vect M$, що при $B=\mathrm{const}$ дає
% ~~~~~~~~~~~~~~~~~~~~~~~~~~~~~~~~~~~~~~~~~~~~~~~~~~~~~~~~
\begin{equation}\label{eq:diamagnetic-current}
	\vect j_\perp=\frac{c}{B^2}\,\vect B\times\nabla p.
\end{equation}
% ~~~~~~~~~~~~~~~~~~~~~~~~~~~~~~~~~~~~~~~~~~~~~~~~~~~~~~~~~~~~~~
Той самий результат одержується з~умови рівноваги $\nabla p=\vect j\times\vect B/c$
(див.~\eqref{eq:equilibrium}). Щоб побачити наслідок, розглянемо плоский шар, у~якому $p=p(x)$ і~поле
$\vect B=B(x)\vect e_z$. Тоді $j_y=-(c/4\pi)\,dB/dx$, і~умова рівноваги $dp/dx=j_yB/c$ дає
% ~~~~~~~~~~~~~~~~~~~~~~~~~~~~~~~~~~~~~~~~~~~~~~~~~~~~~~~~
\begin{equation}\label{eq:pressure-balance}
	p+\frac{B^2}{8\pi}=\mathrm{const}.
\end{equation}
% ~~~~~~~~~~~~~~~~~~~~~~~~~~~~~~~~~~~~~~~~~~~~~~~~~~~~~~~~~~~~~
Всередині плазми поле послаблене, зовні воно сильніше: плазма \enquote{виштовхує} поле,
а~поле стискає плазму. Рівняння~\eqref{eq:pressure-balance} формулює баланс
\emph{газокінетичного} ($p$) і~\emph{магнітного} ($B^2/8\pi$) тисків.\footnote{Строго
кажучи, у~класичній системі в~термодинамічній рівновазі повна намагніченість дорівнює нулю
(теорема Бора~--- ван Левен): діамагнетизм орбіт компенсується струмами частинок, що
ковзають поблизу стінок. Але в~неоднорідній плазмі об'ємні струми $c\,\nabla\times\vect M$
не компенсуються~--- саме вони й~визначають рівновагу.}


%% --------------------------------------------------------
\section{Вмороженість магнітного поля}\index{Вмороженість магнітного поля|textbf}
%% --------------------------------------------------------


З~закону Ома~\eqref{eq:Ohm_law_plasma}

% ========================================================
\begin{equation*}
	 \vect E=-\vect v\times\vect B/c+c\,\rot\vect
	B/(4\pi\sigma).
\end{equation*}
% ========================================================

Підставляючи в~закон індукції $\partial_t\vect B=-c\,\rot\vect E$ і~враховуючи
$\nabla\cdot\vect B=0$, одержуємо \emphx{рівняння індукції}\index{Індукції рівняння|textbf}
% ~~~~~~~~~~~~~~~~~~~~~~~~~~~~~~~~~~~~~~~~~~~~~~~~~~~~~~~~
\begin{equation}\label{eq:induction}
	\frac{\partial\vect B}{\partial t}=\rot(\vect v\times\vect B)+\frac{c^2}{4\pi\sigma}\Delta\vect B.
\end{equation}
% ~~~~~~~~~~~~~~~~~~~~~~~~~~~~~~~~~~~~~~~~~~~~~~~~~~~~~~~~~~~~~~
Перший доданок описує перенесення поля течією, другий~--- дифузію поля. Їхнє відношення
визначає \emphx{магнітне число Рейнольдса}\index{Рейнольдса магнітне число|textbf}
% ~~~~~~~~~~~~~~~~~~~~~~~~~~~~~~~~~~~~~~~~~~~~~~~~~~~~~~~~
\begin{equation*}
	R_m\sim\frac{4\pi\sigma Lv}{c^2}.
\end{equation*}
Для рідких металів у~лабораторії (ртуть, галій) $R_m$ від $10^{-2}$ до~$1$; для токамака $R_m\sim10^{8}$; для сонячної корони $\sim10^{12}$.
Отже, у~космосі вмороженість~--- майже точний закон, а~в~лабораторних рідких металах~--- ні.
% ~~~~~~~~~~~~~~~~~~~~~~~~~~~~~~~~~~~~~~~~~~~~~~~~~~~~~~~~~


\subsection*{Ідеальна провідність (\texorpdfstring{$R_m\gg1$}{R_m >> 1})}


Нехай у~плазмі існує магнітне поле~--- створене зовнішніми котушками, струмом самої плазми
або успадковане від джерела разом із плазмою. Якщо провідність дуже велика, дифузійним
доданком у~\eqref{eq:induction} можна знехтувати, і~рівняння індукції зводиться до
% ~~~~~~~~~~~~~~~~~~~~~~~~~~~~~~~~~~~~~~~~~~~~~~~~~~~~~~~~
\begin{equation*}
    \frac{\partial\vect B}{\partial t}=\rot(\vect v\times\vect B).
\end{equation*}
% ~~~~~~~~~~~~~~~~~~~~~~~~~~~~~~~~~~~~~~~~~~~~~~~~~~~~~~~~
З~нього випливає теорема Альвена: магнітний потік крізь будь-який контур, що рухається
разом із~плазмою, зберігається. Кажуть, що силові лінії \enquote{вморожені} в~плазму~---
плазма не перетинає їх, а~поле рухається разом із~нею. Саме ця властивість лежить в~основі
магнітного утримання: рухаючи плазму, ми переносимо разом із~нею і~поле.

Теорема Альвена не каже, звідки взялося поле: вона описує лише його рух, коли поле вже є. На практиці порядок зазвичай
зворотний~--- \emph{спершу поле, потім плазма}. У~токамаку чи стелараторі котушки вмикають до пуску (вакуум і~нейтральний газ поле не
екранують), а~потім газ іонізують у~готовому полі. У~космосі плазма народжується з~газу, де вже є слабке поле, а~стискання
зберігає потік, і~поле зростає ($B\propto\rho^{2/3}$ при ізотропному стисненні). Так, імовірно, виникає поле радіопульсарів
($\sim10^{12}$~Гс); для магнетарів (від $10^{14}$ до~$10^{15}$~Гс) самого стиснення, найпевніше, замало, і~вважають, що додатково діє динамо.

\emph{Доведення \enquote{на пальцях}.} Потік $\Phi$ крізь контур, що рухається з~плазмою, змінюється з~двох причин:
(а)~через $\partial\vect B/\partial t$ і~(б)~через те, що бічна поверхня, яку описує контур за час $dt$, перетинає силові лінії.
Другий внесок дорівнює $-\oint(\vect v\times\vect B)\cdot d\vect l$. З~$\partial_t\vect B=-c\,\rot\vect E$ і~теореми Стокса
\[\frac{d\Phi}{dt}=-c\oint\Bigl(\vect E+\frac1c\vect v\times\vect B\Bigr)\cdot d\vect l=0,\]
бо $\vect E+\vect v\times\vect B/c=0$ (ідеальний закон Ома). Наслідок: $d(\vect B/\rho)/dt=(\vect B/\rho\cdot\nabla)\vect v$.

Звідси прості закони стискання, які випливають зі збереження потоку $\Phi=BS=\mathrm{const}$. При стисненні лише поперек поля площа
перерізу $S\propto1/\rho$, тому $B\propto\rho$, тобто $B/\rho=\mathrm{const}$ ($B^2\propto\rho^2$, ефективний показник адіабати $\gamma_\perp=2$).
Вздовж поля площа $S$ не змінюється, тому $B=\mathrm{const}$. При ізотропному стисненні всі розміри змінюються однаково: довжина $\ell$,
$S\propto\ell^2$, $\rho\propto\ell^{-3}$, звідки $B\propto1/S\propto\ell^{-2}\propto\rho^{2/3}$. Тоді тиск поля $B^2/8\pi\propto\rho^{4/3}$,
тобто $\gamma=4/3$, як у~випромінювання.


\subsection*{Скінченна провідність: дифузія поля}


У~реальній плазмі провідність скінченна, тому другий доданок у~\eqref{eq:induction} уже
не~можна відкинути. Нехай для простоти плазма нерухома ($\vect v=0$). Тоді рівняння
індукції набирає вигляду рівняння дифузії
% ~~~~~~~~~~~~~~~~~~~~~~~~~~~~~~~~~~~~~~~~~~~~~~~~~~~~~~~~
\begin{equation*}
    \frac{\partial\vect B}{\partial t}
    =\frac{c^2}{4\pi\sigma}\,\Delta\vect B,
\end{equation*}
% ~~~~~~~~~~~~~~~~~~~~~~~~~~~~~~~~~~~~~~~~~~~~~~~~~~~~~~~~
тобто наявне неоднорідне поле з~часом \emphx{розмивається}: струми, що підтримують
неоднорідність, загасають через джоулеві втрати. Характерний час цього розмиття
на~масштабі~$L$
% ~~~~~~~~~~~~~~~~~~~~~~~~~~~~~~~~~~~~~~~~~~~~~~~~~~~~~~~~
\begin{equation*}
    \tau_d=\frac{4\pi\sigma L^2}{c^2}.
\end{equation*}
% ~~~~~~~~~~~~~~~~~~~~~~~~~~~~~~~~~~~~~~~~~~~~~~~~~~~~~~~~
\begin{Estimate}
Лабораторна плазма ($T_e=1$~кеВ, $L=1$~м): $\tau_d\approx50$~с, тому поле треба весь час підтримувати зовнішніми
джерелами. Сонячна корона ($T\sim100$~еВ, $L\sim10^{10}$~см): $\tau_d\sim10^{16}$~с~$\approx3\times10^{8}$~років, що набагато більше за
часи коронних процесів (від хвилин до діб). Тому там дифузією можна знехтувати й~користуватися наближенням вмороженості.
\end{Estimate}

Дифузія діє в~обидва боки: поле розтікається від місця, де його більше, до місця, де менше. Поле, яке було в~плазмі, повільно виходить
назовні; зовнішнє поле так само повільно проникає в~плазму, що його не мала, за час $\tau_d$. Тому \enquote{вморожене} й~\enquote{не проникає} зовні~---
дві сторони однієї властивості: за $R_m\gg1$ поле не заходить у~плазму й~не залишає її, а~лише рухається разом із~нею. Наприклад,
магнітне поле сонячного вітру не проникає в~магнітосферу Землі, а~потрапляє туди лише через перез'єднання (розд.~\ref{sec:reconnection}).


\subsection*{Динамо: як поле підтримується самим рухом}\label{sec:dynamo}


\emphx{Динамо}\index{Динамо|textbf}~--- самопідсилення магнітного поля рухом електропровідного середовища без зовнішнього джерела.
Воно потрібне, бо дифузія розпускає поле за час $\tau_d$: у~рідкому залізному ядрі Землі це $10^4$--$10^5$~років, набагато менше за~її вік
(4,5~млрд років), тож поле має постійно відновлюватися.
Механізм~--- перший доданок рівняння індукції~\eqref{eq:induction}: течія розтягує, скручує й~складає силові лінії, а~розтягнута трубка
потоку за збереження потоку має сильніше поле, як у~законі $B/\rho=\mathrm{const}$ (рис.~\ref{fig:dynamo-stf}).
Потрібні $R_m\gg1$ та~несиметрична течія: осесиметричне поле не можна підтримати осесиметричною течією (теорема Каулінга), тому діють обертання
й~конвекція; на~Сонці диференціальне обертання витягує полоїдальне поле в~тороїдальне ($\Omega$-ефект), а~закручена конвекція знову породжує полоїдальне ($\alpha$-ефект).
Енергію поле бере з~кінетичної енергії течії, і~воно росте, доки сила Лоренца не починає гальмувати рух; так діють земне й~сонячне
(11-річний цикл) динамо, і~так, імовірно, посилюється поле магнетарів.

\begin{figure}[ht]\centering
	\localinput{dynamo-stf.tikz}
	\caption{Схема \enquote{розтягни~--- скрути~--- склади}: петля вмороженого потоку розтягується, скручується у~вісімку й~складається
	у~дві петлі. Потік $\Phi$ зберігається, тому після складання поле зростає вдвічі; повторення циклу підсилює поле експоненційно.}
	\label{fig:dynamo-stf}
\end{figure}

%% --------------------------------------------------------
\section{Перез'єднання силових ліній}\index{Перез'єднання силових ліній|textbf}\label{sec:reconnection}
%% --------------------------------------------------------

В~ідеальній МГД силові лінії вморожені в~плазму, тож топологія поля не~змінюється. Але в~тонкому струмовому шарі,
де поле різко змінює знак, градієнт $B$ великий, і~дифузійний доданок у~\eqref{eq:induction} уже не~можна відкинути (або стають важливими холлівський доданок
та~інерція електронів). Там вмороженість порушується, і~лінії \emphx{перез'єднуються}.

\emph{На~пальцях.} Уявіть дві натягнуті гумки, спрямовані в~протилежні боки; плазма тримає їх на~місці. Їх зближують, і~у~вузькій зоні контакту вмороженість \enquote{вимикається}:
замість двох ліній, що минали одна одну, утворюються дві нові, обернені на~кут (рис.~\ref{pic:reconnection}). Нові лінії різко зігнуті, а~натяг $B^2/4\pi$ вздовж лінії
випрямляє їх, як відпущену тятиву, і~викидає плазму вздовж шару зі~швидкістю порядку альфвенівської. Магнітна енергія перетворюється на~теплову, кінетичну енергію потоків
і~енергію прискорених частинок.

\begin{figure}[ht]\centering
	\localinput{reconnection.tikz}
	\caption{Перез'єднання: (а)~протилежно спрямовані лінії стискають струмовий шар; (б)~після перез'єднання виникає X-точка, а~плазма викидається вздовж шару}
	\label{pic:reconnection}
\end{figure}

\emph{Оцінка швидкості.}\index{Світа--Паркера модель|textbf} У~моделі Світа--Паркера шар довжиною $L$ і~товщиною $\delta$: вхідний потік $v_{\rm in}$ і~вихідний $\sim v_A$ зв'язані збереженням маси,
$v_{\rm in}L\sim v_A\delta$, а~товщина визначається дифузією, $\delta\sim\sqrt{(c^2/4\pi\sigma)\,L/v_A}$. Звідси
% ~~~~~~~~~~~~~~~~~~~~~~~~~~~~~~~~~~~~~~~~~~~~~~~~~~~~~~~~
\begin{equation}\label{eq:sweet-parker}
	\tcbhighmath{\frac{v_{\rm in}}{v_A}\sim R_{m}^{-1/2},\qquad R_m=\frac{4\pi\sigma Lv_A}{c^2}.}
\end{equation}
% ~~~~~~~~~~~~~~~~~~~~~~~~~~~~~~~~~~~~~~~~~~~~~~~~~~~~~~~~
\begin{Estimate}
Для сонячної корони $R_m\sim10^{12}$, тож $v_{\rm in}/v_A\sim10^{-6}$, а~час перез'єднання $\tau\sim(L/v_A)\sqrt{R_m}$ при $L=10^{10}$~см,
$v_A=7\times10^{7}$~см/с сягає $\sim10^{8}$~с (роки), тоді як спалахи тривають хвилини.
\end{Estimate}

Реальне перез'єднання швидше: шар при великому $R_m$ розпадається на~острівці (плазмоїдна нестійкість),
а~у~безсутикових шарах працюють холлівські ефекти (на~масштабі $c/\omega_{pi}$, див.~\eqref{eq:generalized-ohm-from-two-fluid}); швидкість виходить на~$v_{\rm in}$ від $0{,}01$ до~$0{,}1\,v_A$, незалежно від $R_m$.

\begin{Attention}
	Силові лінії не~є матеріальними об'єктами. Мова \enquote{розриву й~зшивання} лише образно описує зміну зв'язності поля: насправді змінюється те,
	які елементи плазми пов'язані однією лінією.
\end{Attention}

\emph{Де спостерігається.} Сонячні спалахи, суббурі в~магнітосфері Землі, пилкоподібні релаксації та~розривні моди в~токамаках (розд.~\ref{sec:instability},~\ref{sec:magnetic-confinement}).

\emph{Проблема нагріву корони.} Корона має $T\sim10^6$~К, у~сотні разів більше за~фотосферу ($\sim6\times10^3$~К), хоча енергія тече від холоднішої поверхні. Потрібне джерело
$\sim10^{5}$--$10^{7}$~ерг/(см$^2\cdot$с). Єдиної відповіді досі немає, але картина така. Енергію несе рух фотосфери, який скручує й~заплітає силові лінії.
Є два шляхи її розсіяння: (1)~\emph{нанофлеєри}~--- безліч малих перез'єднань у~тонких струмових шарах, на~які розпадається заплетене поле (модель Паркера);
(2)~\emph{хвилі}~--- альфвенівська турбулентність, яку запускають рухи вузлів поля у~фотосфері. Огляд Клімчука~\cite{klimchuk2015key} обґрунтовує, що нагрів корони імпульсний, а~корона заповнена
тонкими струмовими шарами, що мають перез'єднуватися; моделювання альфвенівської турбулентності~\cite{vanballegooijen2017heating} дає температуру $\approx2{,}5$~МК
у~петлі активної області, але не~відтворює всіх спостережних обмежень. Нині вважають, що обидва механізми пов'язані: турбулентність створює струмові шари, у~яких і~відбувається перез'єднання.


%% --------------------------------------------------------
\section{Дві мови опису: окремі частинки й~суцільне середовище}\label{sec:two-descriptions}
%% --------------------------------------------------------


Ми описали плазму двічі. В~розд.~\ref{sec:particle-motion} вона~--- набір частинок, кожна
з~яких крутиться навколо своєї силової лінії. Вище~--- суцільне провідне середовище
з~тиском, густиною й~швидкістю течії. Чи не суперечать ці картини одна одній? Здавалося б, суперечать. Візьмемо установку керованого синтезу.

\begin{Estimate}
Для $n=10^{14}~\mathrm{см}^{-3}$, $T=10$~кеВ маємо $\nu_{ei}\approx4{,}4\times10^{3}~\mathrm{с}^{-1}$ і~довжину вільного пробігу
\[
	\lambda_{\mathrm{пр}}\sim\frac{v_{Te}}{\nu_{ei}}\approx\frac{4\times10^{9}~\text{см/с}}{4{,}4\times10^{3}~\mathrm{с}^{-1}}\approx10~\text{км},
\]
у~десять тисяч разів більшу за~розмір установки ($\sim1$~м).
\end{Estimate}

Зіткнень практично немає, а~ми користуємося рівняннями рідини. Але ж рідина, як ми звикли, тримається саме на
зіткненнях.


\subsection*{Що насправді дозволяє перейти до середовища}


У~зіткнювальному газі локальна рівновага підтримується зіткненнями. У~замагніченій плазмі
поперечний рух додатково обмежений гірацією, тому усереднення за ларморівським обертанням
може бути осмисленим навіть за великої довжини вільного пробігу. Але це \emph{не заміна
зіткнень} взагалі: уздовж~$\vect B$ частинки вільно переміщуються, а~розподіл за
швидкостями може залишатися далеким від локальної рівноваги. Саме тому однорідинна МГД має
обмеження.

\begin{Attention}
	Звідси ключове правило. Якщо $L\gg r_{Li}$ та $\omega\ll\omega_{ci}$, макроскопічну
	динаміку часто можна описувати МГД, але лише за наявності відповідного замикання.
	Уздовж~$\vect B$, при анізотропному тиску та в~явищах, чутливих до розподілу частинок
	за швидкостями (загасання Ландау, захоплені частинки в~торі, конус втрат дзеркала),
	однорідинної МГД вже недостатньо й~потрібна кінетика.
\end{Attention}


\subsection*{Словник: що в~частинковій мові відповідає величині рідини}


Подальші три підрозділи показують, як \enquote{наживо} пов'язані ці описи.

\begin{table}[ht]
	\centering\small
	\caption{Словник між частинковим і~гідродинамічним описом.}
	\label{tab:dictionary}
	\begin{platbl}{colspec={X[l,m]X[l,m]X[l,m]}}
		Частинки & Середовище & Зв'язок \\
		Дрейф $\vect E\times\vect B$ однаковий для всіх частинок & швидкість течії $\vect v$ & $\vect E+\vect v\times\vect B/c=0$ \\
		$\sum_sn_sm_s$ & масова густина $\rho_m$ & \\
		$\sum_sn_sT_s$ & тиск $p$ & \\
		Гірація й~градієнт густини & діамагнітний струм $\vect j_\perp$ & $\vect j_\perp=c\,\vect B\times\nabla p/B^2$ \\
		Сума рівнянь руху електронів та іонів & рівняння МГД & $\rho_m\,d\vect v/dt=-\nabla p+\vect j\times\vect B/c$ \\
	\end{platbl}
\end{table}


\subsection*{Вмороженість як наслідок дрейфу}


Закон Ома ідеальної МГД, $\vect E+\vect v\times\vect B/c=0$, не є окремим постулатом.
Дрейф~\eqref{eq:ExB} однаковий для електронів і~іонів, тож уся плазма поперек поля тече
зі~швидкістю $\vect v_E=c\vect E\times\vect B/B^2$. Підставимо $\vect v=\vect v_E$: $\vect
v_E\times\vect B/c=-\vect E_\perp$, тобто $\vect E_\perp+\vect v\times\vect B/c=0$.
А~умова $E_\parallel=0$ означає, що електрони, які вільно рухаються вздовж ліній, миттєво
замикають будь-яке поздовжнє поле. Отже, \enquote{лінії рухаються разом з~плазмою} означає
лише, що частинки прив'язані до ліній, а~самі лінії дрейфують зі швидкістю $\vect v_E$.


\subsection*{Діамагнітний струм \enquote{на пальцях}}


Тут легко помилитися: центри ларморівських кіл не рухаються, а~струм є. Звідки він? Нехай
$\vect B=B\vect e_z$ і~концентрація залежить від~$x$. Розглянемо позитивні частинки
у~вертикальній площині $x=x_0$. З~рівнянь руху $x_g=x+v_y/\omega_c$, де $x_g$~---
координата центра кола. Тобто частинки, що проходять через площину \emph{вгору} ($v_y>0$),
мають центри \emph{праворуч}, а~частинки, що рухаються \emph{вниз},~--- ліворуч. Якщо
справа концентрація більша, вгору летить більше частинок, ніж вниз, і~потоки не
компенсуються:
% ~~~~~~~~~~~~~~~~~~~~~~~~~~~~~~~~~~~~~~~~~~~~~~~~~~~~~~~~
\begin{equation*}
	\Gamma_y=\int v_y\,n\!\left(x_0+\frac{v_y}{\omega_c}\right)f_M\,d^3v
	\simeq\frac{n'}{\omega_c}\langle v_y^2\rangle
	=\frac{n'\,T}{m\,\omega_c}
	=\frac{c\,T}{qB}\,n'.
\end{equation*}
% ~~~~~~~~~~~~~~~~~~~~~~~~~~~~~~~~~~~~~~~~~~~~~~~~~~~~~~~~
Тут $f_M$~--- нормований максвеллівський розподіл за~швидкостями, $n'=dn/dx$, а~нульовий
доданок розкладу зникає через симетрію. Густина струму $j_y=q\Gamma_y=c\,T\,n'/B$. Для електронів і~напрям обертання, і~знак заряду протилежні, тож струми
\emph{додаються}: $j_y=c\,p'/B$, де $p=n(T_e+T_i)$. У~векторному
вигляді $\vect j_\perp=c\,\vect B\times\nabla p/B^2$~--- рівно
\eqref{eq:diamagnetic-current}.

\begin{figure}[ht]\centering
	\localinput{diamagnetic-flux.tikz}
	\caption{Походження діамагнітного струму: центри орбіт нерухомі, але потоки вгору
	й~вниз не компенсуються}
	\label{fig:diamagnetic-flux}
\end{figure}

Той самий струм ми отримали ще двічі: з~намагніченості $M=-p/B$ (розд.~\ref{sec:diamagnetism})
та~з~МГД-рівноваги~\eqref{eq:equilibrium}. Три різних шляхи~--- частинковий,
\enquote{магнітний} і~рідинний~--- дали один результат. Саме так ми й~впевнюємося, що два
описи сумісні.


\subsection*{Діамагнітні дрейфи в~плазмовому циліндрі}


Для утримання важливо бачити, як цей струм виглядає для кожного сорту окремо. Візьмемо
довгий циліндр із~полем $\vect B=B\vect e_z$ і~тиском $p_s(r)$, максимальним на~осі. У~рівновазі
сила тиску, електрична й~магнітна сили в~рідині сорту $s$ (заряд $q_s$) компенсуються:
$q_sn(\vect E+\vect v_s\times\vect B/c)=\nabla p_s$. Векторно множачи на~$\vect B$, дістаємо
% ~~~~~~~~~~~~~~~~~~~~~~~~~~~~~~~~~~~~~~~~~~~~~~~~~~~~~~~~
\begin{equation}\label{eq:diamagnetic-drift}
	\vect v_{\perp s}=\underbrace{c\,\frac{\vect E\times\vect B}{B^2}}_{\vect v_E}
	+\underbrace{\frac{c}{q_sn B^2}\,\vect B\times\nabla p_s}_{\vect v_{*s}}.
\end{equation}
% ~~~~~~~~~~~~~~~~~~~~~~~~~~~~~~~~~~~~~~~~~~~~~~~~~~~~~~~~~~~~~~
Перший доданок~--- дрейф~\eqref{eq:ExB}, однаковий для всіх сортів. Другий називають
\emphx{діамагнітним дрейфом}\index{Діамагнітний дрейф|textbf}: у~нього~$q_s$ стоїть
у~знаменнику, тому він \emph{протилежний} для іонів і~електронів. Це не дрейф
ведучих центрів (вони нерухомі, див.~вище), а~швидкість рідини, зумовлена тим, що
через площину проходить різна кількість частинок. У~циліндрі $\nabla p_s=p_s'\,\vect e_r$ і~$\vect e_z\times\vect e_r=\vect e_\theta$, тому
$$
v_{*i\theta}=\frac{c\,p_i'}{enB},\qquad v_{*e\theta}=-\frac{c\,p_e'}{enB},\qquad
j_\theta=en\,(v_{*i\theta}-v_{*e\theta})=\frac{c\,p'}{B},
$$
де $p=p_i+p_e$, а~вираз для $j_\theta$ збігається з~\eqref{eq:diamagnetic-current}.
Оскільки $p'<0$, іони обертаються навколо осі за~годинниковою стрілкою (якщо дивитися
назустріч~$\vect B$), електрони~--- проти, а~струм тече так, що його поле \emph{послаблює}
$B$ всередині (рис.~\ref{fig:diamagnetic-cylinder}).

\begin{figure}[ht]\centering
	\localinput{diamagnetic-cylinder.tikz}
	\caption{Діамагнітні дрейфи в~циліндрі з~тиском, що спадає від осі. Іони й~електрони
	обертаються в~протилежні боки; сумарний струм послаблює $B$ всередині шнура}
	\label{fig:diamagnetic-cylinder}
\end{figure}

Для утримання з~цього випливає три речі.
\begin{itemize}
	\item \emph{Утримання є балансом сил.} Підставивши $j_\theta$ в~$\nabla p=\vect j\times\vect B/c$, повертаємося
	до~\eqref{eq:pressure-balance}: діамагнітний струм створює саме ту силу $\vect j\times\vect B/c$,
	яка радіально тримає тиск. Чим крутіший профіль, тим більший струм і~тим більше
	послаблюється поле всередині.
	\item \emph{Дрейф замкнений.} У~циліндрі $\vect v_*$ напрямлений уздовж $\theta$, тобто нікуди не веде:
	жодна частинка не покидає шнур, плазма лише обертається як ціле. У~торі ж до~цього додається
	вертикальний дрейф ($\nabla B$ і~кривизна, розд.~\ref{sec:magnetic-confinement}), що розділяє заряди~--- звідси
	потреба в~закрутці ліній поля.
	\item \emph{Швидкість обертання задає радіальне поле.} Сумарна швидкість іонів $v_{i\theta}=v_E+v_{*i}$:
	радіальне поле $E_r$ може вповні скомпенсувати діамагнітний дрейф іонів, і~тоді вони
	спокійні, а~електрони (за $T_e\approx T_i$) обертаються з~подвоєною швидкістю. Зсув $v_\theta(r)$ гасить турбулентні вихори, тому на~ньому
	ґрунтується утворення транспортних бар'єрів (H-мода).
\end{itemize}

\begin{Estimate}
Для $T_i=10$~кеВ, $B=5\times10^4$~Гс і~характерної довжини профілю $L_p=|p/p'|=50$~см
діамагнітна швидкість іонів $v_{*i}=cT_i/(eBL_p)\approx4\times10^5$~см/с~$=4$~км/с. Вона
менша за~теплову $v_{Ti}\sim7\times10^7$~см/с у~$\sim r_{Li}/L_p\sim10^{-2}$ разів, тому
плазма \enquote{тече} повільно, хоча кожна частинка рухається швидко. Цю саму швидкість
$v_*$ має дрейфова хвиля (розд.~\ref{sec:drift-waves}).
\end{Estimate}


\subsection*{Рівняння руху плазми як цілого й~куди зникла електрична сила}


Рівняння руху плазми як цілого ми вивели в~розд.~\ref{sec:mhd}: електронне та іонне рівняння
додаються, тертя скорочується ($\vect R_e+\vect R_i=0$: зіткнення зберігають сумарний
імпульс), і~лишається
% ~~~~~~~~~~~~~~~~~~~~~~~~~~~~~~~~~~~~~~~~~~~~~~~~~~~~~~~~
\begin{equation*}
	\rho_m\frac{d\vect v}{dt}=-\nabla p+\rho_q\vect E+\frac1c\,\vect j\times\vect B.
\end{equation*}
% ~~~~~~~~~~~~~~~~~~~~~~~~~~~~~~~~~~~~~~~~~~~~~~~~~~~~~~~~
Тепер можна пояснити, чому ми раніше відкинули $\rho_q\vect E$. За~Гауссом
$\rho_q=\nabla\cdot\vect E/4\pi\sim E/4\pi L$, тому $\rho_qE\sim E^2/4\pi L$. Натомість
$j\times B/c\sim B^2/4\pi L$. Відношення $E^2/B^2\sim v^2/c^2\ll1$ (бо $E\sim vB/c$):
електричною силою можна знехтувати для повільних ($v\ll c$) рухів~--- тієї ж умови, що
й~для струму зміщення.


\subsection*{Де МГД чесно спрощує}


У~рівнянні руху простої МГД тиск задається скаляром. Це природне наближення для достатньо
частих зіткнень, коли вони встигають ізотропізувати розподіл. У~сильно замагніченій
беззіткнювальній плазмі збереження першого адіабатичного інваріанта приводить до
анізотропії: $p_\perp$ і~$p_\parallel$ можуть відрізнятися. Тоді потрібні кінетичний або
щонайменше двотисковий опис; проста однорідинна МГД вже не містить усієї фізики. Водночас
для довгохвильових мод із майже ізотропним розподілом МГД може залишатися корисним навіть
за слабких зіткнень~--- її застосовність визначається не самим фактом $\lambda_{\mathrm{пр}}>L$,
а~замиканням, частотами та відношенням характерних масштабів до $r_L$ і~$\lambda_D$.


\subsection*{Що це змінює для керованого синтезу}


Кожен із~описів \enquote{відповідає} за свою частину проблеми
(табл.~\ref{tab:fusion-link}).

{\small
\begin{longtblr}[caption={Які ефекти й~моделі що дають для керованого синтезу.},label={tab:fusion-link}]{colspec={X[3,l,m]X[3,l,m]X[5,l,m]}}
		Ефект & Модель & Наслідок для КТС \\
		Дебаївський шар біля стінки й~дивертора & шар Бома (розд.~\ref{sec:bohm-sheath}): $u\ge c_s$, $\varphi_f\approx-3T_e/e$ & усе падіння потенціалу зосереджене у~шарі завтовшки $\sim\lambda_D$; він визначає потік тепла $\sim n c_s(\ldots)T_e$ на матеріал \\
		Відсікання при $\omega<\omega_p$ & $\varepsilon(\omega)=1-\omega_p^2/\omega^2$ & хвиля з~$f_{ce}=\omega_{ce}/(2\pi)\approx1{,}4\times10^{11}$~Гц
для $B=5$~Тл має $n_{\mathrm{кр}}\approx2{,}5\times10^{14}>10^{14}$, тобто $O$-мода проходить у~плазму (розд.~\ref{sec:mhd-waves}) й~нагріває її на циклотронному резонансі \\
		Провідність Спітцера $\sigma\propto T^{3/2}$ & кулонівські зіткнення & омічний нагрів слабшає з~температурою, потрібен додатковий нагрів: пучки нейтралів, ВЧ-хвилі й~$\alpha$-частинки \\
		Дрейфи та збереження $\mu$ & рух окремої частинки & чисто тороїдальне поле не утримує, потрібна закрутка ліній (токамак, стеларатор); дзеркало має конус втрат \\
		$\beta=8\pi p/B^2$ & МГД-рівновага & потужність синтезу $P_f\propto\beta^2B^4$ (див.~нижче), потрібні сильні надпровідні магніти, а~$\beta$ обмежене стійкістю \\
		Класичний і~аномальний перенос & кінетика, турбулентність & класична оцінка дає $\tau\sim10^{3}$~с, реально лише секунди; через це установки доводиться робити великими \\
	\end{longtblr}}

Два рядки таблиці варто виписати. По-перше, чому \enquote{$P_f\propto\beta^2B^4$}. Повна
концентрація частинок $2n$, тому $p=2nT$, і~потужність синтезу
% ~~~~~~~~~~~~~~~~~~~~~~~~~~~~~~~~~~~~~~~~~~~~~~~~~~~~~~~~
\begin{equation*}
	P_f=\tfrac14\,n^2\langle\sigma v\rangle E_f=\frac{p^2}{16\,T^2}\,\langle\sigma v\rangle E_f.
\end{equation*}
% ~~~~~~~~~~~~~~~~~~~~~~~~~~~~~~~~~~~~~~~~~~~~~~~~~~~~~~~~~~~~~~
В~околі $10$--$20$~кеВ у~цьому вузькому інтервалі наведені значення $\langle\sigma
v\rangle$ приблизно зростають як $T^2$. Лише як локальну оцінку в~цьому інтервалі можна
тому писати $P_f\propto p^2=(\beta B^2/8\pi)^2\propto\beta^2B^4$. Це не є універсальним
законом реакції для всіх температур. Подвоєння поля при тому самому~$\beta$ дає у~$16$
разів більшу потужність, і~це головна причина прагнення до сильних полів. Для нашого
прикладу $p=3{,}2\times10^{6}$~дин/см$^2$ ($3{,}2$~бар), магнітний тиск при
$5\times10^4$~Гс дорівнює $\approx100$~бар, тож $\beta\approx3\,\%$.

По-друге, \emph{класичний} перенос. Поперечну теплопровідність іонів оцінюємо як
$\chi_\perp\sim r_{Li}^2\nu_{ii}$ (крок~$r_L$, частота кроків~$\nu_{ii}$). Для
$r_{Li}\approx0{,}3$~см і~$\nu_{ii}\approx\nu_{ei}/(40\text{--}60)\approx10^{2}~\mathrm{с}^{-1}$ маємо
$\chi_\perp\sim10~\mathrm{см}^2/\mathrm{с}$, а~час витоку тепла з~шнура розміру $L\sim1$~м
$\tau\sim L^2/\chi_\perp\sim10^{3}$~с. Тобто по~класиці плазма трималась би сотні секунд
(десятки хвилин). Насправді $\tau_E$ становить секунди~--- на три порядки менше. Різницю дає
\emph{аномальний} перенос, спричинений турбулентністю, і~її не можна отримати ні з~одно-,
ні з~дворідинної МГД: потрібен кінетичний (турбулентний) опис. Це оцінка за~порядком величини: точні
коефіцієнти змінюють результат у~1--10 разів, але не на три порядки.


%% --------------------------------------------------------
\section{Як це вимірюють: магнітні виміри}\label{sec:magnetic-diagnostics}
%% --------------------------------------------------------

У~рідинній картині плазма~--- провідне середовище зі~струмом і~тиском, а~магнітне поле є
її відгуком: струм створює поле за~законом Ампера~$\oint\vect B\cdot d\vect l=4\pi I/c$, а~тиск
послаблює поле всередині за~\eqref{eq:pressure-balance}. Поза плазмою поле можна виміряти
звичайною котушкою. Згідно із~законом Фарадея, котушка з~$N$ витків площі $S$ у~полі $B(t)$
дає ЕРС $\varepsilon=-(NS/c)\,dB/dt$; сигнал інтегрують, і~$\int\varepsilon\,dt\propto B$.
Різні геометрії котушок дають різні величини: струм, енергію, положення шнура.

\subsection*{Пояс Роговського: повний струм}\index{Діагностика плазми!магнітні виміри}\index{Роговського пояс|textbf}
Гнучку котушку з~рівномірним намотуванням ($n_c$ витків на~одиницю довжини, переріз $S$)
замикають навколо плазмового шнура (рис.~\ref{pic:magnetic-diagnostics}\,а). Потік крізь неї
$\Phi=n_cS\oint B_\theta\,dl=(4\pi n_cS/c)\,I$ за~законом Ампера, тому
% ~~~~~~~~~~~~~~~~~~~~~~~~~~~~~~~~~~~~~~~~~~~~~~~~~~~~~~~~
\begin{equation}\label{eq:rogowski}
	\tcbhighmath{\int\varepsilon\,dt=-\frac{4\pi n_cS}{c^2}\,I}\ .
\end{equation}
% ~~~~~~~~~~~~~~~~~~~~~~~~~~~~~~~~~~~~~~~~~~~~~~~~~~~~~~~~
Результат не залежить ні від розподілу струму по~перерізу, ні від форми контуру. Це робить
пояс основним датчиком струму плазми.

\begin{Estimate}
Для $n_c=10$~см$^{-1}$, $S=1$~см$^2$ і~наростання струму $1$~МА за~$0{,}1$~с
ЕРС дорівнює $4\pi n_cS\,\dot I/c^2\approx4\times10^{-3}$~статВ~$\approx1{,}3$~В
($1$~статВ~$\approx300$~В).
\end{Estimate}

\subsection*{Діамагнітна петля: енергія плазми}\index{Діамагнітна петля|textbf}
Петля, що охоплює плазмовий шнур у~площині, перпендикулярній до~$\vect B_0$ (там же, де і~пояс),
вимірює потік поздовжнього поля. З~балансу тисків~\eqref{eq:pressure-balance} при $\beta\ll1$
поле всередині плазми $B_z\approx B_0-4\pi p/B_0$ (рис.~\ref{pic:magnetic-diagnostics}\,б), тобто
потік зменшується на
% ~~~~~~~~~~~~~~~~~~~~~~~~~~~~~~~~~~~~~~~~~~~~~~~~~~~~~~~~
\begin{equation}\label{eq:diamagnetic-loop}
	\tcbhighmath{\Delta\Phi=-\frac{4\pi}{B_0}\int p_\perp\,dS}\ .
\end{equation}
% ~~~~~~~~~~~~~~~~~~~~~~~~~~~~~~~~~~~~~~~~~~~~~~~~~~~~~~~~
Для ізотропної плазми у~тороїдальній системі з~великим радіусом $R$ повна теплова енергія
$W=\tfrac32\int p\,dV\approx\tfrac32\,(2\pi R)\int p\,dS=\tfrac34\,RB_0|\Delta\Phi|$.

\begin{Estimate}
Для $R=150$~см, $a=50$~см, $B_0=5\times10^4$~Гс і~$p=3{,}2\times10^6$~дин/см$^2$
(задача~\ref{prb:PlasmaBeta}) маємо $\int p\,dS=p\pi a^2$ і~$|\Delta\Phi|=4\pi p\,\pi a^2/B_0
\approx6\times10^{6}$~Мкс ($0{,}06$~Вб), тобто лише $\beta/2\approx1{,}6\,\%$ повного
тороїдального потоку $B_0\pi a^2$. Звідси $W\approx3{,}5$~МДж.
\end{Estimate}

\begin{figure}[ht]\centering
	\begin{subfigure}[t]{0.5\linewidth}\centering
		\localinput{magnetic-loops.tikz}
		\caption{}
	\end{subfigure}%
	\begin{subfigure}[t]{0.5\linewidth}\centering
		\localinput{diamagnetic-profile.tikz}
		\caption{}
	\end{subfigure}
	\caption{Магнітні виміри: (а)~переріз плазмового шнура з~поясом Роговського (струм $I_p$), діамагнітною
	петлею (потік $B_t$) і~магнітним зондом на~стінці; (б)~плазма послаблює поле $B_z$ всередині
	($\Delta\Phi\propto\int p\,dS$; ефект перебільшено, вертикальну вісь обрізано)}
	\label{pic:magnetic-diagnostics}
\end{figure}

\subsection*{Магнітні зонди}\index{Магнітні зонди|textbf}
Невелика котушка біля стінки вимірює локальну похідну $\dot B$. Для збурення
$\widetilde B\,e^{-i\omega t}$ сигнал $|\varepsilon|=(NS/c)\,\omega\widetilde B$ зростає з~частотою,
тому зонди чутливі до~швидких збурень. Приклад: $N=100$, $S=1$~см$^2$, $\widetilde B=1$~Гс
і~$f=10$~кГц дають $\approx60$~мВ. За~масивом зондів уздовж обходу шнура визначають
також його положення й~форму (розд.~\ref{sec:beta-limits}).

\subsection*{Межі застосовності}
\begin{itemize}
\item Усі котушки вимірюють лише \emph{зміни} потоку: інтегратор дрейфує, тому тривалі розряди
  (хвилини) потребують особливої точності. Крім того, котушка вимірює все, що створює потік:
  пояс Роговського дає повний струм, включно з~вихровими струмами в~камері.
\item Сигнал діамагнітної петлі становить лише 1--2\,\% потоку тороїдального поля. Його треба
  відокремити від зміни $B_0$ другою, компенсаційною петлею, і~точність результату обмежена
  точністю цього віднімання.
\item Петля вимірює $p_\perp$, а~не повний тиск; у~токамаку на~нього накладається
  парамагнітний внесок полоїдального поля струму, який віднімають за~незалежними вимірами $I_p$.
\item Формула~\eqref{eq:diamagnetic-loop} виведена при $\beta\ll1$ й~у~наближенні балансу
  тисків; для сильно деформованого поля потрібна повна МГД-рівновага.
\item Металева камера екранує швидкі збурення: час проникнення поля крізь стінку обмежує
  смугу частот, тому датчики розташовують усередині камери або враховують стінку.
\end{itemize}

\looseness=-1
Отже, магнітні виміри безпосередньо застосовують три результати цього розділу: закон Ампера
дає струм, баланс тисків~--- енергію, а~закон Фарадея перетворює змінне поле на~електричний
сигнал. Це найпростіші діагностики плазми: вони не потребують лазера, не збурюють плазму й~працюють
на~кожній установці.

%% --------------------------------------------------------
\begin{Summary}
\begin{itemize}
\item Рідинний опис оперує усередненими полями $n$, $\vect v$, $p$ і~не є синонімом частих зіткнень, але потребує замикання.
\item Сума рівнянь руху електронів та іонів дає однорідинну МГД; ідеальна МГД визначається рівнянням $\vect E+\vect v\times\vect B/c=0$.
\item Магнітне поле \enquote{вморожене} в~плазму (вмороженість порушується при $R_m\lesssim1$ і~в~тонких струмових шарах, розд.~\ref{sec:reconnection}); параметр $\beta=8\pi p/B^2$ порівнює газокінетичний і~магнітний тиски.
\item Діамагнітний струм виникає завдяки градієнту тиску й~гірації, хоча центри орбіт нерухомі.
\item Альвенівська хвиля з~швидкістю $v_A=B/\sqrt{4\pi\rho_m}$~--- перша колективна мода замагніченої плазми.
\item Магнітні виміри ґрунтуються на~законах Ампера й~Фарадея та~балансі тисків: пояс Роговського дає струм, діамагнітна петля~--- енергію, зонди~--- швидкі збурення (розд.~\ref{sec:magnetic-diagnostics}).
\end{itemize}
\end{Summary}

\Needspace{9\baselineskip}
\tocunnum{Задачі}
\section*{Задачі}
%% --------------------------------------------------------

\begin{problem}\label{prb:SolenoidPressure}\lvl{2}
	Довгим соленоїдом з~$n$ витками на одиницю довжини тече струм $I$. Знайдіть тиск на бічну поверхню
	соленоїда.
	\begin{solution}
		$p=B^2/8\pi=2\pi n^2I^2/c^2$, $B=4\pi nI/c$.
	\end{solution}
\end{problem}

\begin{problem}\label{prb:PlasmaBeta}\lvl{1}
	Плазма з~$n_e=n_i=10^{14}~\mathrm{см}^{-3}$, $T_e=T_i=10$~кеВ. Знайдіть тиск і~поле,
	при якому $\beta=1$. Яким буде $\beta$ при $B=5\times10^4$~Гс?
	\begin{solution}
		$p=2nT=3{,}2\times10^{6}$~дин/см$^2$, $B(\beta=1)\approx9\times10^{3}$~Гс; $\beta\approx3{,}2\,\%$ при $5\times10^4$~Гс.
	\end{solution}
\end{problem}

\begin{problem}\label{prb:PlasmaDiffusion}\lvl{2}
	Оцініть час дифузії магнітного поля для плазми з~$T_e=1$~кеВ,
	$\ln\Lambda=15$, $L=1$~м. Знайдіть магнітне число Рейнольдса при швидкості руху плазми
	$10^5$~см/с.
	\begin{solution}
		$\sigma\approx3{,}6\times10^{17}~\mathrm{с}^{-1}$ (порівнянна з~міддю), $\tau_d=4\pi\sigma L^2/c^2\approx50$~с, $R_m\approx5\times10^4$.
	\end{solution}
\end{problem}

\begin{problem}\label{prb:PlasmaAlfven}\lvl{1}
	Знайдіть швидкість Альвена для водневої плазми з~$n_i=10^{14}~\mathrm{см}^{-3}$ у~полі
	$B=10^4$~Гс і~порівняйте її з~тепловою швидкістю іонів $v_{Ti}$ при $T_i=10$~кеВ.
	\begin{solution}
		$v_A\approx2{,}2\times10^{8}$~см/с, $v_{Ti}\approx1{,}0\times10^{8}$~см/с ($v_{Ti}/v_A\approx0{,}45$).
	\end{solution}
\end{problem}

\begin{problem}\label{prb:MaxwellStress}\lvl{3}
	Покажіть, що сила на одиницю об'єму $\vect f=\vect j\times\vect B/c$ записується як
	$f_i=\partial_j\sigma_{ij}$ з~тензором~\eqref{eq:maxwell-stress}. Знайдіть
	магнітний тиск при $B=5{,}3$~Тл (у~барах) і~порівняйте з~газовим тиском плазми
	в~задачі~\ref{prb:PlasmaBeta}. Чому дорівнює $\beta$?
	\begin{solution}
		$B^2/8\pi\approx1{,}1\times10^{8}$~дин/см$^2\approx112$~бар, $\beta\approx3\,\%$.
	\end{solution}
\end{problem}

\begin{problem}\label{prb:AlfvenString}\lvl{2}
	Виведіть швидкість альвенівської хвилі, розглядаючи силову лінію як струну з~натягом
	$B^2/4\pi$. Знайдіть $v_A$ для дейтерієвої плазми $n_i=10^{14}~\mathrm{см}^{-3}$,
	$B=5\times10^4$~Гс і~час, за~який хвиля проходить $L=1$~м.
	\begin{solution}
		$v_A\approx7{,}7\times10^{8}$~см/с, $t=L/v_A\approx1{,}3\times10^{-7}$~с.
	\end{solution}
\end{problem}

\begin{problem}\label{prb:FluidScales}\lvl{2}
	Для дейтерієвої плазми з~$n=10^{14}~\mathrm{см}^{-3}$, $T=10$~кеВ, $B=5\times10^4$~Гс оцініть
	$\lambda_{\mathrm{пр}}=v_{Te}/\nu_{ei}$, $r_{Li}$, відношення $D_\perp/D_\parallel$.
	Поясніть, чому МГД застосовна поперек поля, хоча довжина пробігу набагато більша
	за~розмір установки.
	\begin{solution}
		$\lambda_{\mathrm{пр}}\sim10^{6}$~см (10~км), $r_{Li}\approx0{,}3$~см, для іонів $D_\perp/D_\parallel\sim(\nu_{ii}/\omega_{ci})^2\sim10^{-13}$, для електронів $(\nu_{ei}/\omega_{ce})^2\sim10^{-17}$.
	Вздовж поля частинки вільно пролітають $\lambda_{\mathrm{пр}}\gg L$, а~поперек поля гірація з~$r_L\ll L$ локалізує їх: роль довжини пробігу відіграє $r_L$.
	\end{solution}
\end{problem}

\begin{problem}\label{prb:DiamagneticFlux}\lvl{3}
	Для позитивних частинок у~полі $\vect B=B\vect e_z$ з~концентрацією $n(x)$ покажіть,
	що центр кола має координату $x_g=x+v_y/\omega_c$. Знайдіть потік $\Gamma_y$ через
	площину $x=x_0$ і~густину струму $j_y$. Перевірте збіг
	із~\eqref{eq:diamagnetic-current}.
	\begin{solution}
		$\Gamma_y=c\,T\,n^\prime/(qB)$, $j_y=c\,p^\prime/B$.
	\end{solution}
\end{problem}

\begin{problem}\label{prb:FluxTube}\lvl{3}
	Покажіть, що при поперечному стисненні вмороженої силової трубки $B/\rho=\text{const}$
	(довжина трубки не змінюється), а~при ізотропному стисненні $B\propto\rho^{2/3}$.
	У~скільки разів зросте $B$ хмари, що колапсує в~зорю з~ростом густини в~$10^{6}$ разів?
	\begin{solution}
		у~$10^4$ разів.
	\end{solution}
\end{problem}

\begin{problem}\label{prb:ParkerSpiral}\lvl{3}
	Сонце обертається з~$\Omega\approx2{,}7\times10^{-6}$~рад/с, сонячний вітер має швидкість
	$400$~км/с. З~умови вмороженості знайдіть кут між $\vect B$ і~радіальним напрямом на орбіті Землі
	($r=1{,}5\times10^{8}$~км): $\tan\psi=\Omega r/v$.
	\begin{solution}
		$\psi\approx45^\circ$.
	\end{solution}
\end{problem}

\begin{problem}\label{prb:DiamagneticLoop}\lvl{2}
Діамагнітна петля зареєструвала зменшення потоку тороїдального поля $|\Delta\Phi|=3\times10^{6}$~Мкс у~токамаку з~$R=150$~см,
$B_0=5\times10^4$~Гс. Знайдіть теплову енергію плазми $W$, вважаючи тиск ізотропним. Чи можна за~цим сигналом визначити
окремо $n$ і~$T$?
	\begin{solution}
		$W=\tfrac34RB_0|\Delta\Phi|\approx1{,}7\times10^{13}$~ерг~$=1{,}7$~МДж; ні, петля дає лише $\int p\,dS$, тобто добуток~$nT$.
	\end{solution}
\end{problem}

\begin{problem}\label{prb:Rogowski}\lvl{2}
Пояс Роговського має $n_c=20$~см$^{-1}$ і~$S=0{,}5$~см$^2$. Яким буде інтеграл $\int\varepsilon\,dt$ (в~В$\cdot$с) при $I=1$~МА? Чи залежить
результат від положення шнура всередині пояса?
	\begin{solution}
		$\mu_0n_cSI$ в~одиницях СІ $=4\pi\!\times\!10^{-7}\cdot2000\cdot0{,}5\!\times\!10^{-4}\cdot10^{6}\approx0{,}13$~В$\cdot$с;
не залежить~--- лише від повного охопленого струму.
	\end{solution}
\end{problem}

%% --------------------------------------------------------

\begin{Questions}
\begin{quizlist}
\item Чому плазма діамагнітна? Який зв'язок між діамагнетизмом
		  і~рівнянням~\eqref{eq:pressure-balance}?

\item Що таке магнітний тиск і~натяг силових ліній? Як вони входять в~рівняння
		  МГД?

\item Що означає вмороженість поля? Від якого параметра залежить, чи вона
		  виконується?

\item Як із~закону Ампера й~сили Лоренца отримати тензор натягів і~розкласти силу
		  на тиск і~натяг? Чому магнітна сила ніколи не діє вздовж поля?

\item Чому плазма майже без зіткнень описується рівняннями рідини поперек поля,
		  але не вздовж? Яку роль виконує гірація?

\item Чому вмороженість поля є наслідком дрейфу $\vect E\times\vect B$? Звідки
		  береться діамагнітний струм, якщо центри орбіт нерухомі?

\item Чому джоулеве тепло не порушує збереження повної енергії?

\item Коли МГД непридатна (анізотропія тиску, беззіткнювальна плазма вздовж $\vect B$)?

\item Чому силові лінії поводяться як натягнуті гумові нитки, що відштовхуються одна від одної?

\item Чому діамагнітна петля вимірює енергію плазми, а~пояс Роговського~--- струм? Які закони лежать в~основі кожного з~них?

\end{quizlist}
\end{Questions}

\Closesolutionfile{answer}
